\documentclass[12pt,a4paper]{report}
\usepackage[utf8]{inputenc}
\usepackage[brazil]{babel}
\usepackage{amsmath}
\usepackage{amsfonts}
\usepackage{amssymb}
\usepackage{enumitem}
\usepackage{fancyhdr}
\usepackage{float}
\usepackage{geometry}
\usepackage{graphicx}
\usepackage{multicol}
\usepackage{tcolorbox}
\usepackage{parskip}
\usepackage{setspace}
\usepackage{tikz}
\usepackage{pgfplots}
\pgfplotsset{compat=1.15}
\usepackage{verbatim}
\usetikzlibrary{arrows.meta, calc}
\setstretch{1.2}
\tcbuselibrary{skins, breakable}
\geometry{top=2cm, bottom=2cm, left=2.5cm, right=2.5cm}
\usepackage[colorlinks=true, linkcolor=blue, urlcolor=blue, citecolor=blue]{hyperref}

\usepackage{chngcntr}
\counterwithout{footnote}{chapter}

\pagestyle{fancy}
\fancyhf{} 
\fancyhead[C]{\leftmark} 
\fancyfoot[C]{\thepage} 

% Estilo de caixa para destaque conceitual
\newtcolorbox{destaque}[1][]{
	enhanced,
	breakable,
	colback=gray!5,
	colframe=black!70,
	boxrule=1pt,
	arc=4pt,
	left=10pt,
	right=10pt,
	top=8pt,
	bottom=8pt,
	fonttitle=\bfseries,
	code={%
		\ifstrempty{#1}{%
			\pgfkeysalso{notitle}%
		}{%
			\pgfkeysalso{title={#1}}%
		}%
	}
}

\begin{document}
	
	\begin{titlepage}
		\centering
		
		\includegraphics[width=8cm]{../alexjordane.github.io/ifes.png}\\[1cm]
		
		{\Large \textbf{Instituto Federal do Espírito Santo}}\\[0.5cm]
		{\large \textbf{Licenciatura em Matemática}}\\[0.5cm]
		{\large Introdução ao Estudo das Funções}
		
		\vspace{\fill}
		
		{\Huge \textbf{Logaritmos e Função Logarítmica}}\\[0.5cm]
		{\Large Material de Apoio}
		
		\vspace{\fill}
		
		{\Large Professor Alex Jordane}
		
		\vspace{\fill}
		
		{\large Vitória -- ES}\\[0.2cm]
		{\large 2026}
		
	\end{titlepage}
	
	\clearpage
	\setcounter{page}{1}  % Inicia a contagem pré-textual a partir de i
	
	\chapter*{Apresentação}
	\addcontentsline{toc}{chapter}{Apresentação}
	
	O estudo dos logaritmos e da função logarítmica constitui um momento singular na formação matemática. Longe de se reduzir a um conjunto de regras de manipulação simbólica ou a artifícios técnicos de simplificação, a criação dos logaritmos representa uma resposta histórica às necessidades de cálculo de sua época e uma ponte conceitual indispensável para a compreensão de fenômenos contínuos.
	
	Nossa investigação parte de uma situação-problema real ligada à farmacocinética clínica. Diante do decaimento da concentração de um medicamento na corrente sanguínea, deparamo-nos com uma relação exponencial cujas bases distintas não admitem redução direta por fatoração elementar. É exatamente essa necessidade prática que motiva a emergência dos logaritmos como instrumento capaz de converter potências e variações exponenciais em relações lineares operáveis.
	
	Ao longo dos capítulos, articulamos o conceito de logaritmo, suas propriedades operatórias, a relação de inversa da função exponencial, o comportamento gráfico e as técnicas de resolução de equações e inequações. Esperamos que este texto sirva como um recurso aberto e acolhedor de consulta, reflexão e estudo autônomo, apoiando estudantes e docentes na construção de compreensões sólidas acerca do pensamento funcional.
	
	Bons estudos!
	
	\vspace{\fill}
	
	{\footnotesize
		\texttt{Esta é uma primeira versão do material, criada em 02/11/2026.}\\
		\texttt{Este material foi elaborado em colaboração com a inteligência artificial Gemini, \\ versão 3.1 Pro.} \\		
		\texttt{O compartilhamento e a reprodução deste material são permitidos, desde que a fonte \\ e a autoria sejam citadas.}  \\		
		\texttt{Esta obra está licenciada com uma Licença Atribuição-NãoComercial-SemDerivações 4.0.}
		\vspace{1cm}
		\begin{center}
			\includegraphics[width=4cm]{CC.png}
	\end{center}}
	
	\clearpage
	
	\chapter*{Convenções e Notações Matemáticas}
	\addcontentsline{toc}{chapter}{Convenções e Notações Matemáticas}
	\thispagestyle{empty}
	\label{chap:convencoes}
	
	Para facilitar a leitura e manter a clareza, reunimos as principais notações e convenções adotadas neste material:
	
	\begin{itemize}[label=$\circ$]
		\item \textbf{Conjuntos numéricos e restrições:}
		\begin{itemize}[label=--]
			\item $\mathbb{N} = \{0, 1, 2, \dots\}$, $\mathbb{Z}$, $\mathbb{Q}$ e $\mathbb{R}$.
			\item $\mathbb{R}^*_+ = \{x \in \mathbb{R} \mid x > 0\} = {]0, +\infty[}$ (domínio fundamental dos logaritmos).
		\end{itemize}
		
		\item \textbf{Intervalos reais (notação com colchetes invertidos):}
		\begin{itemize}[label=--]
			\item $[a, b]$: fechado em ambos os extremos ($\{x \in \mathbb{R} \mid a \leqslant x \leqslant b\}$).
			\item $]a, b[$: aberto em ambos os extremos ($\{x \in \mathbb{R} \mid a < x < b\}$).
			\item $]a, +\infty[$ e $]-\infty, b[$: abertos e ilimitados.
		\end{itemize}
		
		\item \textbf{Logaritmos:}
		\begin{itemize}[label=--]
			\item $\log_a b$: logaritmo de $b$ na base $a$ ($a > 0$, $a \neq 1$ e $b > 0$).
			\item $\log x = \log_{10} x$ (base decimal omitida).
			\item $\ln x = \log_e x$ (logaritmo natural, com $e \approx 2{,}718$).
			\item $\text{antilog}_a x = a^x$.
		\end{itemize}
		
		\item \textbf{Pares ordenados e separadores:}
		\begin{itemize}[label=--]
			\item Em coordenadas com vírgula decimal, utiliza-se ponto e vírgula como separador: $P(x_0; y_0)$ (exemplo: $(3{,}5; 0)$).
			\item Adota-se a vírgula para números decimais (padrão brasileiro) e o símbolo $\blacksquare$ para indicar o encerramento de demonstrações.
		\end{itemize}
	\end{itemize}	
	\clearpage
	
	\tableofcontents
	\thispagestyle{empty}
	
	\clearpage
	%\pagenumbering{arabic} % Retorna à numeração arábica tradicional (1, 2, 3, ...)
	%\setcounter{page}{1}   % Reinicia a contagem para que o Capítulo 1 comece na página 1
		
	\chapter{Uma situação problema real}
	\label{chap:problema}
	
	Vamos começar esse material com um problema muito comum em nosso cotidiano. A \textbf{farmacocinética clínica} se ocupa em determinar o tempo necessário para que a concentração de um medicamento caia a um patamar seguro ou terapêutico no organismo. Nosso problema está relacionado a este contexto.
	
	\begin{destaque}[O Problema do Ajuste Medicamentoso]
		\label{dest:problema}
		Em um tratamento hospitalar, uma pessoa recebe uma dose de determinado analgésico por via venosa. A partir do momento em que a substância entra na corrente sanguínea, o organismo passa a eliminá-la de forma contínua por meio dos processos metabólicos.
		
		Estudos mostram que a cada hora a concentração do medicamento diminui em $20\%$, ou seja, resta $80\%$ ($0{,}8$) da quantidade da hora anterior.
		
		A concentração inicial no sangue é de $50\text{ mg/L}$. Para que o medicamento continue aliviando a dor sem atingir níveis de toxicidade, a equipe de saúde precisa saber em quantas horas a concentração atingirá exatamente $15\text{ mg/L}$, momento planejado para a aplicação de uma nova dose ou avaliação clínica.
		
		Sabendo isso, como calculamos o momento exato que deve ser aplicado nova dose?
	\end{destaque}
	
	\section{O modelo matemático}
	A concentração $C(t)$ em função do tempo decorrido $t$ (em horas) é descrita pela relação exponencial decrescente:
	$$C(t) = 50 \cdot (0{,}8)^t$$
	Igualando a concentração ao valor desejado de $15\text{ mg/L}$:
	$$50 \cdot (0{,}8)^t = 15$$
	Dividindo ambos os lados por $50$:
	$$(0{,}8)^t = \frac{15}{50}$$
	$$(0{,}8)^t = 0{,}3$$
	Na forma fracionária irredutível:
	$$\left(\frac{4}{5}\right)^t = \frac{3}{10}$$
	
	\section{O problema da redução a bases comuns}
	
	Aqui entramos em um dilema, como resolver essa equação?	
	
	Até este ponto dos estudos de equações exponenciais, resolvemos equações buscando bases comuns por meio da fatoração simples (como transformar $4^x = 32$ em $2^{2x} = 2^5$). 
	
	No entanto, ao analisar a relação obtida:
	\begin{itemize}[label=$\circ$]
		\item Do lado esquerdo, temos potências com fatores primos $2$ e $5$.
		\item Do lado direito, surge o fator primo $3$ no numerador, além dos fatores $2$ e $5$ no denominador.
		\item Não existe expoente racional simples que permita igualar diretamente as bases $\frac{4}{5}$ e $\frac{3}{10}$ pelo método tradicional de decomposição.
	\end{itemize}
	
	É justamente esse impasse que justifica a introdução de um novo conceito, o de \textbf{logaritmo}. Esse novo conceito será a ferramenta algébrica para extrair o expoente quando a redução a bases comuns não é viável.
	
	Vamos começar a estudar essa nova ferramenta no próximo capítulo e em breve retornaremos ao problema \ref{dest:problema}.
	
	
	\chapter{Logaritmos}
	\label{chap:logaritmos}
	
	\section{Conceito e Definição}
	
	Vimos que equações como $\left(\frac{4}{5}\right)^t = \frac{3}{10}$ nos colocam diante da busca por um expoente real que relaciona duas grandezas positivas quando a fatoração elementar não é suficiente. Formalizamos essa operação inversa da exponenciação por meio do conceito de logaritmo.
	
	\begin{destaque}[Definição de Logaritmo]
		Dados os números reais $a$ e $b$, com $a > 0$, $a \neq 1$ e $b > 0$, chama-se \textbf{logaritmo de $b$ na base $a$} o expoente $x$ ao qual se deve elevar a base $a$ para obter o número $b$:
		$$\log_a b = x \iff a^x = b$$
	\end{destaque}
	
	Nessa relação:
	\begin{itemize}[label=$\circ$]
		\item $a$ é a \textbf{base} do logaritmo;
		\item $b$ é o \textbf{logaritmando} (ou argumento);
		\item $x$ é o \textbf{logaritmo}.
	\end{itemize}
	
	As restrições impostas ($a > 0$, $a \neq 1$ e $b > 0$) constituem as chamadas \textbf{condições de existência}:
	\begin{itemize}[label=$\circ$]
		\item Se $a \leqslant 0$, a potência $a^x$ não estaria bem definida para expoentes reais quaisquer (por exemplo, $(-4)^{1/2} \notin \mathbb{R}$).
		\item Se $a = 1$, teríamos $1^x = 1$ para todo $x$, impossibilitando determinar um valor único para expressões como $\log_1 5$ ou tornando $\log_1 1$ indeterminado.
		\item Como $a > 0$, qualquer potência real de $a$ resulta sempre em um número estritamente positivo; logo, é indispensável que $b > 0$.
	\end{itemize}
	
	\subsection*{Exemplos de Aplicação Direta}
	
	\begin{enumerate}[label=\textbf{\alph*)}]
		\item $\log_2 16 = 4$, pois $2^4 = 16$.
		\item $\log_3 \left(\frac{1}{3}\right) = -1$, pois $3^{-1} = \frac{1}{3}$.
		\item $\log_7 \sqrt{7} = \frac{1}{2}$, pois $7^{1/2} = \sqrt{7}$.
		\item $\log_61 = 0$, pois $6^0 = 1$.
		\item $\log_{1/2} 16 = -4$, pois $\left(\frac{1}{2}\right)^{-4} = 2^4 = 16$.
	\end{enumerate}
	
	\section{Exercícios Resolvidos}
	
	\subsection*{Exercícios Resolvidos}
	
	\textbf{Exercício 1.} Calcule o valor da expressão:
	$$E = \log_2 64 - \log_3 \left(\frac{1}{27}\right) + \log_{10} 0{,}01$$
	
	\textit{Resolução:}
	Calculamos cada termo separadamente a partir da definição:
	\begin{itemize}[label=$\circ$]
		\item $\log_2 64 = x \iff 2^x = 64 \iff 2^x = 2^6 \iff x = 6$.
		\item $\log_3 \left(\frac{1}{27}\right) = y \iff 3^y = \frac{1}{27} \iff 3^y = 3^{-3} \iff y = -3$.
		\item $\log_{10} 0{,}01 = z \iff 10^z = \frac{1}{100} \iff 10^z = 10^{-2} \iff z = -2$.
	\end{itemize}
	Substituindo os valores obtidos na expressão:
	$$E = 6 - (-3) + (-2) = 6 + 3 - 2 = 7$$
	
	\vspace{0.3cm}
		
	\subsection*{Exercícios Propostos}
	
	\begin{enumerate}
		\item Calcule o valor dos seguintes logaritmos:
		\begin{multicols}{2}
			\begin{enumerate}[label=\textbf{\alph*)}]
				\item $\log_4 16$
				\item $\log_3 \sqrt[3]{9}$
				\item $\log_{1/3} 81$
				\item $\log_{0{,}5} 32$
			\end{enumerate}
		\end{multicols}
		
		\item Determine o valor de $x$ em cada uma das sentenças:
		\begin{multicols}{2}
			\begin{enumerate}[label=\textbf{\alph*)}]
				\item $\log_x 49 = 2$
				\item $\log_2 x = -4$
				\item $\log_4 (2x - 1) = 2$
				\item $\log_{x+1} 25 = 2$
			\end{enumerate}
		\end{multicols}
		
	\end{enumerate}
	
	\section{Consequências Imediatas da Definição:}
	
	Da equivalência fundamental $\log_a b = x \iff a^x = b$, decorrem quatro consequências para todo $a > 0$, $a \neq 1$ e termos pertencentes aos seus respectivos domínios:
	
	\begin{destaque}[Consequências]
		\begin{enumerate}[label=\textbf{\arabic*.}]
			\item \textbf{Logaritmo da unidade:} O logaritmo de $1$ em qualquer base é igual a zero.
			$$\log_a 1 = 0 \quad (\text{pois } a^0 = 1)$$
			
			\item \textbf{Logaritmo da própria base:} O logaritmo da base na mesma base é igual a um.
			$$\log_a a = 1 \quad (\text{pois } a^1 = a)$$
			
			\item \textbf{Potência de expoente logarítmico:} A potência de base $a$ e expoente $\log_a b$ é igual a $b$.
			$$a^{\log_a b} = b \quad (\text{para todo } b > 0)$$
				
			\item \textbf{Unicidade do logaritmo:} Dois logaritmos em uma mesma base são iguais se, e somente se, os logaritmandos forem iguais.
			$$\log_a b = \log_a c \iff b = c \quad (\text{para } b, c > 0)$$		
			
		\end{enumerate}
	\end{destaque}
	
	\begin{destaque}[Demonstração da Consequência 3]
		Sejam $a > 0$, $a \neq 1$ e $b > 0$. 
		
		Pela definição fundamental de logaritmo, temos:
		$$\log_a b = x \iff a^x = b$$
		
		Como $x$ representa exatamente a expressão $\log_a b$, substituímos $x$ diretamente no expoente da equação exponencial $a^x = b$:
		$$a^{\log_a b} = b$$
		
		Essa identidade expressa o fato de que a função exponencial de base $a$ e a função logarítmica de mesma base são operações inversas, desfazendo uma à outra.
		$$\blacksquare$$
	\end{destaque}

		
	\begin{destaque}[Demonstração da Consequência 4]
		$$\log_a b = \log_a c \Longleftrightarrow a^{\log_a c} = b \Longleftrightarrow c=b$$
		$$\blacksquare$$
	\end{destaque}
	
	\subsection*{Exemplos Numéricos}
	
	\begin{enumerate}[label=\textbf{\alph*)}]
		\item $\log_8 1 = 0$
		\item $\log_{\pi} \pi = 1$
		\item $3^{\log_3 7} = 7$
		\item $5^{\log_5 12} = 12$
		\item $2^{1 + \log_2 5} = 2^1 \cdot 2^{\log_2 5} = 2 \cdot 5 = 10$
		\item Se $\log_5 (3x - 2) = \log_5 (x + 8)$, temos $3x - 2 = x + 8 \implies 2x = 10 \implies x = 5$.
	\end{enumerate}
	
	\section{Sistemas de Logaritmos}
	
	Embora a base $a$ possa ser qualquer número real positivo e diferente de $1$, dois sistemas assumem papel de relevo em aplicações científicas e tecnológicas:
	
	\subsection{Sistema de Logaritmos Decimais (Base 10)}
	Criado e sistematizado por Henry Briggs (1561--1630), utiliza a base $10$, a mesma do nosso sistema de numeração posicional. Por convenção usual, quando a base é $10$, omite-se o índice na notação:
	$$\log_{10} x = \log x$$
	
	Exemplos:
	\begin{itemize}[label=$\circ$]
		\item $\log 10 = 1$
		\item $\log 1000 = \log (10^3) = 3$
		\item $\log 0{,}001 = \log (10^{-3}) = -3$
	\end{itemize}
	
	\subsection{Sistema de Logaritmos Naturais ou Neperianos (Base $e$)}
	Utiliza como base o número irracional $e \approx 2{,}71828\dots$ (número de Euler). Esse sistema surge espontaneamente no estudo de fenômenos de crescimento e decaimento contínuos na natureza, na física e na economia. 
	
	O logaritmo na base $e$ é denominado \textbf{logaritmo natural} (ou \textbf{neperiano}, em homenagem ao matemarico escocês John Napier (1550-1617)) e representado pela notação $\ln$:
	$$\log_e x = \ln x$$
	
	Exemplos:
	\begin{itemize}[label=$\circ$]
		\item $\ln e = 1$
		\item $\ln 1 = 0$
		\item $\ln (e^4) = 4$
		\item $e^{\ln 5} = 5$
	\end{itemize}
	
	\section{Propriedades Operatórias dos Logaritmos}
	
	As propriedades operatórias constituem a principal razão histórica e analítica do desenvolvimento dos logaritmos, pois permitem transformar operações aritméticas de maior complexidade (multiplicações, divisões e potenciações) em operações mais simples (adições, subtrações e multiplicações por constantes).
	
	\subsection{Logaritmo do Produto}
	
	A primeira propriedade estabelece que o logaritmo de um produto de dois termos estritamente positivos é igual à soma dos logaritmos desses termos tomados na mesma base.
	
	\begin{destaque}[Logaritmo do Produto]
		Sejam $a, b, c \in \mathbb{R}$ tais que $a > 0$, $a \neq 1$, $b > 0$ e $c > 0$. Então:
		$$\log_a (b \cdot c) = \log_a b + \log_a c$$
	\end{destaque}
	
	\begin{destaque}[Demonstração]
		Para demonstrar essa relação, nomeamos os logaritmos individuais de cada fator:
		$$\log_a b = x \quad \text{e} \quad \log_a c = y$$
		
		Pela definição fundamental de logaritmo, reescrevemos essas igualdades na forma exponencial:
		$$a^x = b \quad \text{e} \quad a^y = c$$
		
		Multiplicando membro a membro as duas igualdades obtidas:
		$$b \cdot c = a^x \cdot a^y$$
		
		Pelas propriedades fundamentais das potências de mesma base, conservamos a base e somamos os expoentes:
		$$b \cdot c = a^{x + y}$$
		
		Aplicando novamente a definição de logaritmo a essa relação exponencial:
		$$\log_a (b \cdot c) = x + y$$
		
		Substituindo $x$ e $y$ pelas expressões originais:
		$$\log_a (b \cdot c) = \log_a b + \log_a c$$
		$$\blacksquare$$
	\end{destaque}
	
	\begin{destaque}[Observação: Produto de termos de mesmo sinal]
		A condição de existência do logaritmo exige apenas que o logaritmando seja estritamente positivo, ou seja, $b \cdot c > 0$. Para que essa desigualdade se verifique no conjunto dos números reais, há duas possibilidades:
		\begin{enumerate}[label=\alph*)]
			\item $b > 0$ e $c > 0$ (ambos positivos);
			\item $b < 0$ e $c < 0$ (ambos negativos).
		\end{enumerate}
		
		No segundo caso, as expressões $\log_a b$ e $\log_a c$ isoladas não existem em $\mathbb{R}$, pois os argumentos seriam negativos. No entanto, como seus valores absolutos $|b| = -b > 0$ e $|c| = -c > 0$ são termos positivos, temos que $b \cdot c = (-b) \cdot (-c) = |b| \cdot |c|$. 
		
		Portanto, de maneira abrangente, se $b \cdot c > 0$, a propriedade assume a forma:
		$$\log_a (b \cdot c) = \log_a |b| + \log_a |c|$$
		
		\textbf{Exemplo numérico:}  
		Considerando $b = -4$ e $c = -8$ na base $2$, temos $(-4) \cdot (-8) = 32 > 0$. Logo:
		$$\log_2 \bigl((-4) \cdot (-8)\bigr) = \log_2 32 = 5$$
		Aplicando a relação com módulos:
		$$\log_2 |-4| + \log_2 |-8| = \log_2 4 + \log_2 8 = 2 + 3 = 5$$
	\end{destaque}
	
	\subsection*{Exemplos}
	
	\begin{enumerate}[label=\textbf{\alph*)}]
		\item Sabendo que $\log 2 \approx 0{,}30$ e $\log 3 \approx 0{,}48$, podemos calcular:
		$$\log 6 = \log(2 \cdot 3) = \log 2 + \log 3 \approx 0{,}30 + 0{,}48 = 0{,}78$$
		
		\item $\log_2 (8 \cdot 16) = \log_2 8 + \log_2 16 = 3 + 4 = 7$.
	\end{enumerate}
	
	\subsection*{Extensão para o Produto de $n$ Fatores}
	
	A propriedade do produto se estende naturalmente para qualquer quantidade finita $n \geqslant 2$ de fatores reais estritamente positivos.
	
	\newpage
	
	\begin{destaque}[Logaritmo do Produto de $n$ Fatores]
		Sejam $a > 0$ com $a \neq 1$, e $b_1, b_2, \dots, b_n \in \mathbb{R}^*_+$, com $n \geqslant 2$. Então:
		$$\log_a (b_1 \cdot b_2 \dotsm b_n) = \log_a b_1 + \log_a b_2 + \dots + \log_a b_n$$
		Ou, em notação de produtório e somatório:
		$$\log_a \left(\prod_{i=1}^n b_i\right) = \sum_{i=1}^n \log_a b_i$$
	\end{destaque}
	
	\begin{destaque}[Demonstração por Indução Finita]
		A demonstração é feita por indução matemática sobre o número de fatores $n$.
		
		\textbf{1. Base da indução ($n = 2$):}  
		Para dois fatores, a proposição é verdadeira conforme demonstrado anteriormente:
		$$\log_a (b_1 \cdot b_2) = \log_a b_1 + \log_a b_2$$
		
		\textbf{2. Hipótese de indução:}  
		Assumimos que a propriedade seja válida para $k$ fatores ($k \geqslant 2$), isto é:
		$$\log_a (b_1 \cdot b_2 \dotsm b_k) = \log_a b_1 + \log_a b_2 + \dots + \log_a b_k$$
		
		\textbf{3. Passo indutivo:}  
		Devemos provar a validade da relação para $k + 1$ fatores:
		$$\log_a (b_1 \cdot b_2 \dotsm b_k \cdot b_{k+1})$$
		
		Associando os primeiros $k$ fatores em um único bloco $B = (b_1 \cdot b_2 \dotsm b_k)$, temos o produto de dois termos: $B$ e $b_{k+1}$. Como $b_i > 0$ para todo $i$, segue que $B > 0$. Aplicando o caso base:
		$$\log_a \bigl(B \cdot b_{k+1}\bigr) = \log_a B + \log_a b_{k+1}$$
		
		Pela hipótese de indução, expandimos $\log_a B$:
		$$\log_a B = \log_a (b_1 \cdot b_2 \dotsm b_k) = \log_a b_1 + \log_a b_2 + \dots + \log_a b_k$$
		
		Substituindo essa soma na expressão anterior:
		$$\log_a (b_1 \cdot b_2 \dotsm b_k \cdot b_{k+1}) = \log_a b_1 + \log_a b_2 + \dots + \log_a b_k + \log_a b_{k+1}$$
		
		Pelo princípio da indução matemática, a propriedade é válida para todo número inteiro $n \geqslant 2$.
		$$\blacksquare$$
	\end{destaque}
	
	\subsection{Logaritmo do Quociente}
	
	A segunda propriedade operatória estabelece que o logaritmo da divisão (ou razão) entre dois números estritamente positivos é igual à diferença entre os logaritmos do dividendo (numerador) e do divisor (denominador), mantida a mesma base.
	
	\begin{destaque}[Logaritmo do Quociente]
		Sejam $a, b, c \in \mathbb{R}$ tais que $a > 0$, $a \neq 1$, $b > 0$ e $c > 0$. Então:
		$$\log_a \left(\frac{b}{c}\right) = \log_a b - \log_a c$$
	\end{destaque}
	
	\begin{destaque}[Demonstração]
		Nomeamos os logaritmos do numerador e do denominador:
		$$\log_a b = x \quad \text{e} \quad \log_a c = y$$
		
		Pela definição fundamental de logaritmo, escrevemos essas relações na forma exponencial:
		$$a^x = b \quad \text{e} \quad a^y = c$$
		
		Dividindo membro a membro a primeira igualdade pela segunda:
		$$\frac{b}{c} = \frac{a^x}{a^y}$$
		
		Pelas regras operatórias das potências de mesma base, mantemos a base e subtraímos os expoentes:
		$$\frac{b}{c} = a^{x - y}$$
		
		Aplicando novamente a definição de logaritmo à igualdade exponencial obtida:
		$$\log_a \left(\frac{b}{c}\right) = x - y$$
		
		Substituindo $x$ e $y$ por suas expressões originais:
		$$\log_a \left(\frac{b}{c}\right) = \log_a b - \log_a c$$
		$$\blacksquare$$
	\end{destaque}
	
	\begin{destaque}[Caso Particular: Logaritmo do Inverso]
		Como decorrência imediata, fazendo $b = 1$ e sabendo que $\log_a 1 = 0$, obtemos o logaritmo do inverso de um número:
		$$\log_a \left(\frac{1}{c}\right) = \log_a 1 - \log_a c = 0 - \log_a c = -\log_a c$$
	\end{destaque}
	
	\begin{destaque}[Observação: Quociente de termos de mesmo sinal]
		De modo inteiramente análogo ao produto, a existência de $\log_a \left(\frac{b}{c}\right)$ requer unicamente que a razão seja estritamente positiva, ou seja, $\frac{b}{c} > 0$. Isso ocorre em duas situações:
		\begin{enumerate}[label=\alph*)]
			\item $b > 0$ e $c > 0$ (numerador e denominador positivos);
			\item $b < 0$ e $c < 0$ (numerador e denominador negativos).
		\end{enumerate}
		
		Quando ambos são negativos, as expressões isoladas $\log_a b$ e $\log_a c$ não estão definidas no conjunto dos números reais. Entretanto, como a divisão de dois valores negativos resulta em um número positivo idêntico à divisão de seus respectivos valores absolutos, temos:
		$$\frac{b}{c} = \frac{-b}{-c} = \frac{|b|}{|c|}$$
		
		Dessa forma, caso saibamos apenas que $\frac{b}{c} > 0$, a propriedade assume a formulação abrangente:
		$$\log_a \left(\frac{b}{c}\right) = \log_a |b| - \log_a |c|$$
		
		\textbf{Exemplo numérico:}  
		Tomando $b = -32$ e $c = -4$ na base $2$, verificamos que $\frac{-32}{-4} = 8 > 0$. Assim:
		$$\log_2 \left(\frac{-32}{-4}\right) = \log_2 8 = 3$$
		Utilizando a forma modular:
		$$\log_2 |-32| - \log_2 |-4| = \log_2 32 - \log_2 4 = 5 - 2 = 3$$
	\end{destaque}
	
	\subsection*{Exemplos}
	
	\begin{enumerate}[label=\textbf{\alph*)}]
		\item Sabendo que $\log 2 \approx 0{,}30$ e $\log 3 \approx 0{,}48$, calculemos $\log 1{,}5$:
		$$\log 1{,}5 = \log\left(\frac{3}{2}\right) = \log 3 - \log 2 \approx 0{,}48 - 0{,}30 = 0{,}18$$
		
		\item O cálculo de $\log 5$ em base decimal é uma aplicação clássica da propriedade do quociente, escrevendo $5 = \frac{10}{2}$:
		$$\log 5 = \log\left(\frac{10}{2}\right) = \log 10 - \log 2 = 1 - 0{,}30 = 0{,}70$$
		
		\item $\log_3 \left(\frac{1}{81}\right) = -\log_3 81 = -\log_3 (3^4) = -4$.
		
		\item $\log_2 \left(\frac{32}{4}\right) = \log_2 32 - \log_2 4 = 5 - 2 = 3$.
	\end{enumerate}

	\subsection{Logaritmo de Potência}
	
	A terceira propriedade operatória estabelece que o logaritmo de uma potência cuja base é estritamente positiva equivale ao produto do expoente pelo logaritmo da base dessa potência (conhecida intuitivamente como a ``regra do tombo'').
	
	\begin{destaque}[Logaritmo de uma Potência]
		Sejam $a, b \in \mathbb{R}$ tais que $a > 0$, $a \neq 1$, $b > 0$, e seja $\alpha \in \mathbb{R}$. Então:
		$$\log_a (b^\alpha) = \alpha \cdot \log_a b$$
	\end{destaque}
	
	\begin{destaque}[Demonstração]
		Definimos o logaritmo da base $b$:
		$$\log_a b = x$$
		
		Pela definição fundamental de logaritmo, escrevemos essa relação na forma exponencial:
		$$a^x = b$$
		
		Elevando ambos os membros da igualdade ao expoente real $\alpha$:
		$$(a^x)^\alpha = b^\alpha$$
		
		Pelas propriedades da potenciação (potência de potência):
		$$a^{\alpha \cdot x} = b^\alpha$$
		
		Aplicando a definição de logaritmo a essa nova igualdade exponencial:
		$$\log_a (b^\alpha) = \alpha \cdot x$$
		
		Substituindo $x$ pela sua expressão original ($\log_a b$):
		$$\log_a (b^\alpha) = \alpha \cdot \log_a b$$
		$$\blacksquare$$
	\end{destaque}
	
	\begin{destaque}[Caso Particular: Logaritmo de um Radical]
		Como todo radical pode ser expresso como uma potência de expoente fracionário, fazendo $\alpha = \frac{1}{n}$ (com $n \in \mathbb{N}$ e $n \geqslant 2$), decorre que:
		$$\log_a (\sqrt[n]{b}) = \log_a \left(b^{1/n}\right) = \frac{1}{n} \cdot \log_a b = \frac{\log_a b}{n}$$
	\end{destaque}
	
	\begin{destaque}[Observação: Potências de Expoente Par e Valores Absolutos]
		Se o expoente $\alpha$ for um número inteiro par ($\alpha = 2k$, com $k \in \mathbb{Z}^*$), a potência $b^{2k}$ é estritamente positiva para todo $b \neq 0$, o que assegura a existência de $\log_a (b^{2k})$ mesmo quando $b$ é um número negativo.
		
		Todavia, para que o termo isolado $\log_a b$ faça sentido no conjunto dos números reais, a base do logaritmando precisa ser positiva. Logo, caso saibamos unicamente que $b \neq 0$:
		$$\log_a (b^{2k}) = 2k \cdot \log_a |b|$$
		
		\textbf{Exemplo numérico:}  
		Para $b = -3$ e a base $a = 10$:
		$$\log \bigl((-3)^2\bigr) = \log 9 = \log(3^2) = 2 \cdot \log |-3| = 2 \cdot \log 3$$
		Omissões do módulo conduziriam à expressão indefinida $2 \cdot \log(-3)$.
	\end{destaque}
	
	\subsection*{Exemplos}
	
	\begin{enumerate}[label=\textbf{\alph*)}]
		\item Sabendo que $\log 2 \approx 0{,}30$, calculemos $\log 32$:
		$$\log 32 = \log(2^5) = 5 \cdot \log 2 \approx 5 \cdot 0{,}30 = 1{,}50$$
		
		\item Aplicação em radicais:
		$$\log_3 (\sqrt[5]{27}) = \log_3 \left(27^{1/5}\right) = \frac{1}{5} \cdot \log_3 (3^3) = \frac{3}{5}$$
		
		\item Determinação de logaritmos decimais por potências inversas:
		$$\log \left(\frac{1}{\sqrt{10}}\right) = \log\left(10^{-1/2}\right) = -\frac{1}{2} \cdot \log 10 = -\frac{1}{2} \cdot 1 = -0{,}5$$
	\end{enumerate}
	
	\section{Desenvolvimento e Redução de Expressões Logarítmicas}
	
	A aplicação combinada das propriedades operatórias permite dois procedimentos algébricos complementares:
	\begin{itemize}[label=$\circ$]
		\item \textbf{Desenvolvimento logarítmico:} transformar o logaritmo de uma expressão algébrica complexa (envolvendo multiplicações, divisões, potências e raízes) em uma soma algébrica de logaritmos simples.
		\item \textbf{Redução (ou condensação) a um único logaritmo:} realizar o caminho inverso, reunindo vários logaritmos em um único logaritmando.
	\end{itemize}
	
	\subsection{Desenvolvimento de uma Expressão Algébrica}
	
	Tomemos como modelo analítico a expressão geral:
	$$A = \frac{a^x \cdot \sqrt[y]{b}}{c^z}$$
	
	Assumindo as condições de existência ($a, b, c > 0$, base do logaritmo $k > 0$ com $k \neq 1$, e índice do radical $y \in \mathbb{N}$ com $y \geqslant 2$), aplicamos o logaritmo em ambos os membros:
	$$\log_k A = \log_k \left( \frac{a^x \cdot \sqrt[y]{b}}{c^z} \right)$$
	
	O procedimento de expansão ocorre em etapas sucessivas:
	
	\begin{destaque}[Passo a Passo do Desenvolvimento]
		\begin{enumerate}[label=\textbf{\arabic*.}]
			\item \textbf{Aplicação da propriedade do quociente:}  
			O logaritmo da razão converte-se na diferença entre o logaritmo do numerador e o do denominador:
			$$\log_k A = \log_k \left( a^x \cdot \sqrt[y]{b} \right) - \log_k (c^z)$$
			
			\item \textbf{Aplicação da propriedade do produto:}  
			No primeiro termo, o logaritmo do produto decompõe-se na soma dos logaritmos:
			$$\log_k A = \log_k (a^x) + \log_k (\sqrt[y]{b}) - \log_k (c^z)$$
			
			\item \textbf{Conversão do radical em potência:}  
			Reescrevemos a raiz como expoente fracionário ($\sqrt[y]{b} = b^{1/y}$):
			$$\log_k A = \log_k (a^x) + \log_k \left( b^{1/y} \right) - \log_k (c^z)$$
			
			\item \textbf{Aplicação da propriedade das potências:}  
			Os expoentes passam a multiplicar os respectivos logaritmos:
			$$\log_k A = x \cdot \log_k a + \frac{1}{y} \cdot \log_k b - z \cdot \log_k c$$
		\end{enumerate}
	\end{destaque}
	
	\subsection*{Exemplos Numéricos e Algébricos}
	
	\textbf{Exemplo 1 (Desenvolvimento Numérico).} Desenvolva a expressão $E = \frac{x^3 \cdot \sqrt{y}}{z^4}$ em termos de logaritmos decimais de $x$, $y$ e $z$ (com $x, y, z > 0$):
	$$\log E = \log\left(\frac{x^3 \cdot y^{1/2}}{z^4}\right) = \log(x^3) + \log\left(y^{1/2}\right) - \log(z^4) = 3\log x + \frac{1}{2}\log y - 4\log z$$
	
	\vspace{0.3cm}
	
	\vspace{0.3cm}
	
	\textbf{Exemplo 2 (Processo Inverso: Redução a um Único Logaritmo).}  
	Escreva a expressão $P = 2\log_a x + \frac{1}{3}\log_a y - 3\log_a z$ como o logaritmo de uma única expressão algébrica:
	\begin{align*}
		P &= \log_a (x^2) + \log_a \left(y^{1/3}\right) - \log_a (z^3) \\
		&= \log_a \left(x^2 \cdot \sqrt[3]{y}\right) - \log_a (z^3) \\
		&= \log_a \left(\frac{x^2 \cdot \sqrt[3]{y}}{z^3}\right)
	\end{align*}
	
	\section{Exercícios sobre Propriedades Operatórias}
	
	\subsection{Exercícios Resolvidos}
	
	\textbf{Exercício 1.} Sabendo que $\log 2 \approx 0{,}30$ e $\log 3 \approx 0{,}48$, determine o valor de:
	\begin{enumerate}[label=\textbf{\alph*)}]
		\item $\log 72$
		\item $\log \sqrt{0{,}6}$
	\end{enumerate}
	
	\textit{Resolução:}
	\begin{enumerate}[label=\textbf{\alph*)}]
		\item Decompondo $72$ em fatores primos: $72 = 8 \cdot 9 = 2^3 \cdot 3^2$.  
		Aplicando as propriedades do produto e da potência:
		\begin{align*}
			\log 72 &= \log(2^3 \cdot 3^2) = \log(2^3) + \log(3^2) \\
			&= 3 \cdot \log 2 + 2 \cdot \log 3 \\
			&\approx 3 \cdot (0{,}30) + 2 \cdot (0{,}48) = 0{,}90 + 0{,}96 = 1{,}86
		\end{align*}
		
		\item Escrevemos a raiz e o número decimal em forma fracionária: $\sqrt{0{,}6} = \left(\frac{6}{10}\right)^{1/2} = \left(\frac{3}{5}\right)^{1/2}$.  
		Lembrando que $5 = \frac{10}{2}$, podemos expressar:
		$$\log 5 = \log\left(\frac{10}{2}\right) = \log 10 - \log 2 \approx 1 - 0{,}30 = 0{,}70$$
		Assim:
		\begin{align*}
			\log \sqrt{0{,}6} &= \frac{1}{2} \cdot \log\left(\frac{3}{5}\right) = \frac{1}{2} \cdot [\log 3 - \log 5] \\
			&\approx \frac{1}{2} \cdot [0{,}48 - 0{,}70] = \frac{1}{2} \cdot (-0{,}22) = -0{,}11
		\end{align*}
	\end{enumerate}
	
	\vspace{0.3cm}
	
	\textbf{Exercício 2.}  O potencial hidrogeniônico ($\text{pH}$) mede o nível de acidez ou alcalinidade de uma solução aquosa, sendo definido por:
	$$\text{pH} = -\log [\text{H}^+]$$
	em que $[\text{H}^+]$ representa a concentração de íons hidrogênio em mol por litro ($\text{mol/L}$).  
	Em uma análise de controle da qualidade da água potável, verificou-se que uma amostra continha $[\text{H}^+] = 3{,}2 \times 10^{-7}\text{ mol/L}$. Considerando $\log 2 \approx 0{,}30$, determine o $\text{pH}$ dessa amostra e indique se a água está ligeiramente ácida ($\text{pH} < 7$) ou neutra/alcalina ($\text{pH} \geqslant 7$).
	
	\textit{Resolução:}  
	Substituindo a concentração na relação do $\text{pH}$:
	$$\text{pH} = -\log (3{,}2 \times 10^{-7})$$
	
	Expressamos $3{,}2$ na forma de fração com potências conhecidas: $3{,}2 = \frac{32}{10} = \frac{2^5}{10}$.  
	Substituindo na expressão do produto:
	$$3{,}2 \times 10^{-7} = \frac{2^5}{10} \times 10^{-7} = 2^5 \times 10^{-8}$$
	
	Aplicamos as propriedades do produto e da potência:
	\begin{align*}
		\log (2^5 \times 10^{-8}) &= \log(2^5) + \log(10^{-8}) \\
		&= 5 \cdot \log 2 + (-8) \cdot \log 10 \\
		&= 5 \cdot \log 2 - 8
	\end{align*}
	
	Substituindo a aproximação $\log 2 \approx 0{,}30$:
	$$\log (3{,}2 \times 10^{-7}) \approx 5 \cdot (0{,}30) - 8 = 1{,}50 - 8 = -6{,}50$$
	
	Portanto:
	$$\text{pH} = -(-6{,}50) = 6{,}50$$
	Como $\text{pH} = 6{,}50 < 7$, a água analisada encontra-se ligeiramente ácida.
	
	\vspace{0.3cm}
	
	\textbf{Exercício 3.}  A sensibilidade do ouvido humano a variações de pressão sonora não é linear, mas logarítmica. O nível de intensidade sonora $N$, expresso em decibéis ($\text{dB}$), é dado pela relação:
	$$N = 10 \cdot \log \left(\frac{I}{I_0}\right)$$
	em que $I$ é a intensidade física da onda sonora (em $\text{W/m}^2$) e $I_0 = 10^{-12}\text{ W/m}^2$ é o limiar de audibilidade humana.  
	Em um ambiente escolar, a intensidade do ruído durante o intervalo foi medida em $I_1 = 2 \times 10^{-5}\text{ W/m}^2$, enquanto na sala de aula durante a aula era de $I_2 = 2 \times 10^{-8}\text{ W/m}^2$.  
	Usando as propriedades dos logaritmos, determine:
	\begin{enumerate}[label=\textbf{\alph*)}]
		\item O nível sonoro $N_1$ no intervalo (adote $\log 2 \approx 0{,}30$).
		\item A diferença $N_1 - N_2$ expressa a partir do quociente das intensidades.
	\end{enumerate}
	
	\textit{Resolução:}
	\begin{enumerate}[label=\textbf{\alph*)}]
		\item Calculamos a razão entre $I_1$ e $I_0$:
		$$\frac{I_1}{I_0} = \frac{2 \times 10^{-5}}{10^{-12}} = 2 \times 10^7$$
		Aplicando as propriedades do produto e da potência no logaritmo:
		\begin{align*}
			N_1 &= 10 \cdot \log(2 \times 10^7) = 10 \cdot [\log 2 + \log(10^7)] \\
			&= 10 \cdot [\log 2 + 7 \cdot \log 10] \\
			&\approx 10 \cdot [0{,}30 + 7] = 10 \cdot 7{,}30 = 73\text{ dB}
		\end{align*}
		
		\item A diferença de níveis sonoros é dada por:
		$$N_1 - N_2 = 10 \cdot \log\left(\frac{I_1}{I_0}\right) - 10 \cdot \log\left(\frac{I_2}{I_0}\right) = 10 \cdot \left[ \log\left(\frac{I_1}{I_0}\right) - \log\left(\frac{I_2}{I_0}\right) \right]$$
		Pela propriedade do quociente de logaritmos:
		$$N_1 - N_2 = 10 \cdot \log\left( \frac{I_1 / I_0}{I_2 / I_0} \right) = 10 \cdot \log\left(\frac{I_1}{I_2}\right)$$
		Substituindo a razão $\frac{I_1}{I_2} = \frac{2 \times 10^{-5}}{2 \times 10^{-8}} = 10^3$:
		$$N_1 - N_2 = 10 \cdot \log(10^3) = 10 \cdot 3 = 30\text{ dB}$$
		Esse resultado mostra que um aumento de $1\,000$ vezes na intensidade física da onda corresponde a um acréscimo de apenas $30\text{ dB}$ na percepção auditiva.
	\end{enumerate}
	

	\subsection{Exercícios Propostos}
	
	\begin{enumerate}
		\item Dados $\log 2 \approx 0{,}301$ e $\log 3 \approx 0{,}477$, determine o valor numérico aproximado de:
		\begin{multicols}{2}
			\begin{enumerate}[label=\textbf{\alph*)}]
				\item $\log 18$
				\item $\log 24$
				\item $\log 0{,}75$
				\item $\log \sqrt[3]{12}$
			\end{enumerate}
		\end{multicols}
		
		\item Reduza as seguintes expressões a um único logaritmo:
		\begin{enumerate}[label=\textbf{\alph*)}]
			\item $A = \log_a x + 2\log_a y - 3\log_a z$ \quad ($x, y, z > 0$)
			\item $B = \frac{1}{2}\log_b (x + 1) - \log_b y - 4\log_b z$ \quad ($x > -1$, $y, z > 0$)
		\end{enumerate}
		
		\item (\textbf{Escala Richter}) A magnitude $M$ de um terremoto na escala Richter relaciona-se com a energia liberada $E$ (em joules) por meio da expressão:
		$$\log E = 4{,}4 + 1{,}5 \cdot M$$
		\begin{enumerate}[label=\textbf{\alph*)}]
			\item Sabendo que dois terremotos registraram magnitudes $M_1 = 7{,}0$ e $M_2 = 5{,}0$, use a propriedade do quociente de logaritmos para determinar a razão $\frac{E_1}{E_2}$ entre as energias liberadas.
			\item Interprete o que essa razão representa em termos de liberação de energia física entre dois graus inteiros da escala.
		\end{enumerate}
		
		\item Mostre que, se $a, b \in \mathbb{R}^*_+$ e $a \neq 1$, então:
		$$\log_a \left(\frac{a\cdot (a + b)}{b}\right) = 1 - \log_a \left(\frac{b}{a+b}\right)$$
	\end{enumerate}
	
	\section{Mudança de Base}
	
	Em muitas situações analíticas, deparamo-nos com logaritmos expressos em bases distintas daquelas disponíveis em tabelas, instrumentos computacionais (que operam comumente nas bases $10$ e $e$) ou convenientes para a simplificação algébrica de uma equação. A relação de mudança de base permite reescrever qualquer logaritmo em termos de uma nova base escolhida arbitrariamente.
	
	\begin{destaque}[Fórmula de Mudança de Base]
		Sejam $a, b, c \in \mathbb{R}$ tais que $a > 0$, $a \neq 1$, $b > 0$, $c > 0$ e $c \neq 1$. O logaritmo de $b$ na base $a$ pode ser expresso na nova base $c$ pela razão:
		$$\log_a b = \frac{\log_c b}{\log_c a}$$
	\end{destaque}
	
	\newpage
	
	\begin{destaque}[Demonstração:]
		Definimos inicialmente:
		$$\log_a b = x$$
		
		Pela definição fundamental de logaritmo, escrevemos a relação correspondente na forma exponencial:
		$$a^x = b$$
		
		Como ambos os membros são estritamente positivos e $c$ é uma base admissível ($c > 0$ e $c \neq 1$), aplicamos o logaritmo na base $c$ a ambos os lados da igualdade:
		$$\log_c (a^x) = \log_c b$$
		
		Aplicando a propriedade do logaritmo de uma potência no primeiro membro:
		$$x \cdot \log_c a = \log_c b$$
		
		Como $a \neq 1$, temos que $\log_c a \neq 0$. Logo, podemos isolar $x$:
		$$x = \frac{\log_c b}{\log_c a}$$
		
		Substituindo $x$ por $\log_a b$, obtemos a relação desejada:
		$$\log_a b = \frac{\log_c b}{\log_c a}$$
		$$\blacksquare$$
	\end{destaque}
	
	\newpage
	
	\subsection{Consequências Imediatas}
	
	A fórmula de mudança de base dá origem a identidades práticas de simplificação algébrica:
	
	\begin{destaque}[Consequências Práticas]
		\begin{enumerate}[label=\textbf{\arabic*.}]
			\item \textbf{Inversão de base e logaritmando:}  
			Fazendo $c = b$ na relação fundamental (com $b \neq 1$), temos:
			$$\log_a b = \frac{\log_b b}{\log_b a} = \frac{1}{\log_b a}$$
			Ou seja:
			$$\log_a b \cdot \log_b a = 1$$
			
			\item \textbf{Regra da cadeia para produtos de logaritmos:}  
			Para termos positivos e bases admissíveis $a, b, c$:
			$$\log_a b \cdot \log_b c = \frac{\log_k b}{\log_k a} \cdot \frac{\log_k c}{\log_k b} = \frac{\log_k c}{\log_k a} = \log_a c$$
			De modo geral, para bases e argumentos encadeados:
			$$\log_{a_1} a_2 \cdot \log_{a_2} a_3 \dotsm \log_{a_{n-1}} a_n = \log_{a_1} a_n$$
			
			\item \textbf{Expoente na base:}  
			Para $\beta \in \mathbb{R}^*$:
			$$\log_{a^\beta} b = \frac{\log_a b}{\log_a (a^\beta)} = \frac{\log_a b}{\beta \cdot \log_a a} = \frac{1}{\beta} \cdot \log_a b$$
			Combinando com o expoente no logaritmando:
			$$\log_{a^\beta} (b^\alpha) = \frac{\alpha}{\beta} \cdot \log_a b$$
		\end{enumerate}
	\end{destaque}
	
	\subsection*{Exemplos}
	
	\begin{enumerate}[label=\textbf{\alph*)}]
		\item \textbf{Cálculo com dados em base 10:}  
		Sabendo que $\log 2 \approx 0{,}30$ e $\log 3 \approx 0{,}48$, determine $\log_2 3$:
		$$\log_2 3 = \frac{\log 3}{\log 2} \approx \frac{0{,}48}{0{,}30} = 1{,}60$$
		
		\item \textbf{Simplificação de expoente na base:}  
		$$\log_8 16 = \log_{2^3} (2^4) = \frac{4}{3} \cdot \log_2 2 = \frac{4}{3} \cdot 1 = \frac{4}{3}$$
		
		\item \textbf{Simplificação por encadeamento de bases:}  
		$$\log_2 3 \cdot \log_3 5 \cdot \log_5 8 = \log_2 8 = \log_2 (2^3) = 3$$
	\end{enumerate}
	
	\section{Nota Histórica}
	\begin{destaque}[As Tábuas de Logaritmos e o Cálculo Numérico]
		Antes do surgimento dos recursos computacionais modernos, cálculos astronômicos, náuticos e geodésicos demandavam meses de trabalho manual exaustivo com multiplicações e divisões de números com muitas casas decimais. A invenção dos logaritmos, no início do século XVII, operou uma verdadeira revolução científica ao substituir o esforço dessas operações pela simples adição e subtração.
		
		O matemático escocês John Napier (1550--1617) dedicou mais de duas décadas ao desenvolvimento do método, publicado em 1614 em sua obra \textit{Mirifici Logarithmorum Canonis Descriptio}. Pouco tempo depois, em diálogo com o matemático inglês Henry Briggs (1561--1630), consolidou-se o sistema de base decimal, ideal para o cálculo numérico em virtude de nossa base de numeração posicional. Briggs publicou tábuas com dezenas de milhares de logaritmos calculados com extraordinária precisão (inicialmente com catorze casas decimais).
		
		Na prática, essas tábuas funcionavam como instrumentos essenciais de mediação do trabalho científico:
		\begin{itemize}[label=$\circ$]
			\item Para calcular produtos ou quocientes complexos, determinavam-se os logaritmos dos fatores na tabela;
			\item Realizava-se a soma ou subtração desses valores;
			\item Pelo processo inverso (busca pelo antilogaritmo na tábua), localizava-se o número correspondente ao resultado final.
		\end{itemize}
		
		Quando o valor do logaritmo obtido caía entre duas entradas consecutivas da tábua, recorria-se à interpolação linear para estimar com precisão as casas decimais adicionais. 
		
		O célebre astrônomo e matemático Pierre-Simon Laplace resumiu esse impacto ao afirmar que a criação dos logaritmos, ao encurtar o trabalho dos astrônomos, prolongou-lhes a vida útil dedicada à ciência.
	\end{destaque}
	
	
	\chapter{Retorno ao Ponto de Partida: A Resolução do Ajuste Medicamentoso}
	
	Ao final do Capítulo~\ref{chap:problema}, deparamo-nos com o dilema prático da equipe de saúde: determinar em quantas horas a concentração sanguínea de um analgésico atingiria exatamente $15\text{ mg/L}$, partindo de $50\text{ mg/L}$ e decrescendo a uma taxa de $20\%$ por hora. A modelagem conduziu à equação exponencial:
	$$\left(\frac{4}{5}\right)^t = \frac{3}{10} \quad \text{ou, na forma decimal,} \quad (0{,}8)^t = 0{,}3$$
	
	Naquele momento, a fatoração clássica em bases comuns mostrou-se insuficiente devido à incompatibilidade dos fatores primos dos dois lados da igualdade ($2$ e $5$ à esquerda; $3$, $2$ e $5$ à direita). Munidos agora das propriedades operatórias e dos conceitos de logaritmo, podemos finalmente resolver essa relação com rigor algébrico.
	
	\begin{destaque}[Resolução do Problema Clínico]
		Partimos da equação exponencial:
		$$(0{,}8)^t = 0{,}3$$
		
		Para expressar os valores em termos de dados tabelados usuais, adotamos:
		$$\log 2 \approx 0{,}30 \quad \text{e} \quad \log 3 \approx 0{,}48$$
		
		\begin{enumerate}[label=\textbf{\arabic*.}]
			\item \textbf{Aplicação de logaritmos decimais:}  
			Aplicando o logaritmo na base $10$ em ambos os membros da igualdade:
			$$\log(0{,}8^t) = \log(0{,}3)$$
			\newpage
			\item \textbf{Uso da propriedade da potência:}  
			Pela regra do tombo, o expoente $t$ passa a multiplicar o logaritmo da base:
			$$t \cdot \log(0{,}8) = \log(0{,}3)$$
			Isolando o tempo $t$:
			$$t = \frac{\log(0{,}3)}{\log(0{,}8)}$$
			
			\item \textbf{Decomposição por meio das frações decimais:}  
			Escrevendo os argumentos na forma fracionária decimal:
			$$0{,}3 = \frac{3}{10} \quad \text{e} \quad 0{,}8 = \frac{8}{10} = \frac{2^3}{10}$$
			
			\item \textbf{Cálculo do numerador e do denominador:}  
			Aplicando as propriedades do quociente e da potência, lembrando que $\log 10 = 1$:
			\begin{itemize}[label=$\circ$]
				\item \textbf{Numerador:}  
				$$\log(0{,}3) = \log\left(\frac{3}{10}\right) = \log 3 - \log 10 \approx 0{,}48 - 1 = -0{,}52$$
				\item \textbf{Denominador:}  
				$$\log(0{,}8) = \log\left(\frac{2^3}{10}\right) = 3 \cdot \log 2 - \log 10 \approx 3 \cdot (0{,}30) - 1 = 0{,}90 - 1 = -0{,}10$$
			\end{itemize}
			
			\item \textbf{Determinação do tempo $t$ e interpretação clínica:}  
			Substituindo os valores encontrados:
			$$t = \frac{-0{,}52}{-0{,}10} = 5{,}2\text{ horas}$$
		\end{enumerate}
		
		Para traduzir a fração de hora em minutos de forma precisa para a equipe hospitalar:
		$$0{,}2\text{ h} = 0{,}2 \times 60\text{ minutos} = 12\text{ minutos}$$
		
		Portanto, a concentração atinge exatamente o patamar terapêutico planejado de 15~mg/L após \textbf{5 horas e 12 minutos} da aplicação venosa inicial.
	\end{destaque}
	
	\section*{Síntese}
	
	O percurso percorrido evidencia o papel que conferiu aos logaritmos sua relevância histórica e sua importância analítica:
	\begin{itemize}[label=$\circ$]
		\item Diante de fenômenos contínuos em que as grandezas não coincidem em potências inteiras simples, o logaritmo atua como uma possibilidade que transforma o expoente desconhecido em um fator linear.
		\item A razão entre dois números negativos ($\frac{-0{,}52}{-0{,}10}$) resultou em um tempo estritamente positivo ($5{,}2\text{ horas}$), expressando coerentemente a natureza decrescente do processo de eliminação do fármaco pelo organismo.
	\end{itemize}
	
	Com a base conceitual e operatória dos logaritmos devidamente estabelecida, estamos prontos para avançar, no próximo capítulo, para a formalização da \textbf{Função Logarítmica}, seu comportamento gráfico, domínio e propriedades analíticas.
	
	\chapter{A Função Logarítmica}
	\label{chap:funcao_logaritmica}
	
	Após o estudo dos logaritmos como ferramenta algébrica e operatória, nosso olhar se volta para o comportamento funcional. Vamos analisar como a variação contínua do logaritmando produz variações no valor do logaritmo, caracterizando uma nova classe de funções elementares: as funções logarítmicas.
	
	\section{Definição e Condições de Existência}
	
	\begin{destaque}[Definição de Função Logarítmica]
		Dado um número real $a$, com $a > 0$ e $a \neq 1$, chama-se \textbf{função logarítmica de base $a$} a aplicação $f \colon \mathbb{R}^*_+ \to \mathbb{R}$ que associa a cada número real estritamente positivo $x$ o seu logaritmo na base $a$:
		$$f(x) = \log_a x$$
	\end{destaque}
	
	A partir dessa definição, destacam-se imediatamente os conjuntos fundamentais da função:
	\begin{itemize}[label=$\circ$]
		\item \textbf{Domínio:} $\text{Dom}(f) = \mathbb{R}^*_+ = {]0, +\infty[}$. A função só admite valores reais quando o argumento for estritamente positivo.
		\item \textbf{Contradomínio e Conjunto Imagem:} O contradomínio é $\mathbb{R}$ e a imagem coincide com todo o conjunto dos números reais, isto é, $\text{Im}(f) = \mathbb{R} = {]-\infty, +\infty[}$.
	\end{itemize}
	
	\subsection*{Determinação de Domínio em Expressões Compostas}
	
	Quando trabalhamos com funções em que o logaritmando ou a base dependem de uma variável real (da forma $g(x) = \log_{b(x)} [u(x)]$), o domínio da função é o conjunto de todos os valores de $x$ para os quais as condições de existência sejam simultaneamente satisfeitas:
	$$\begin{cases}
		u(x) > 0 & \text{(o logaritmando deve ser estritamente positivo)} \\
		b(x) > 0 & \text{(a base deve ser estritamente positiva)} \\
		b(x) \neq 1 & \text{(a base deve ser diferente de $1$)}
	\end{cases}$$
	
	\textbf{Exemplo 1.} Determine o domínio da função real $f(x) = \log_2 (3x - 12)$.
	
	\textit{Resolução:}  
	A base é fixa ($2 > 0$ e $2 \neq 1$). A única exigência recai sobre o logaritmando:
	$$3x - 12 > 0 \implies 3x > 12 \implies x > 4$$
	Logo:
	$$\text{Dom}(f) = \{x \in \mathbb{R} \mid x > 4\} = {]4, +\infty[}$$
	
	\vspace{0.3cm}
	
	\textbf{Exemplo 2.} Determine o domínio da função real $g(x) = \log_{(x - 2)} (9 - x^2)$.
	
	\textit{Resolução:}  
	Precisamos impor as três condições fundamentais:
	\begin{enumerate}[label=\textbf{\arabic*.}]
		\item Logaritmando positivo:
		$$9 - x^2 > 0 \implies x^2 < 9 \implies -3 < x < 3 \quad (\text{Intervalo } {]-3, 3[})$$
		\item Base positiva:
		$$x - 2 > 0 \implies x > 2 \quad (\text{Intervalo } {]2, +\infty[})$$
		\item Base diferente de 1:
		$$x - 2 \neq 1 \implies x \neq 3$$
	\end{enumerate}
	
	Determinando a interseção das três condições:
	$${]-3, 3[} \;\cap\; {]2, +\infty[} \;\cap\; (\mathbb{R} - \{3\}) = {]2, 3[}$$
	Portanto, o domínio de $g$ é:
	$$\text{Dom}(g) = \{x \in \mathbb{R} \mid 2 < x < 3\} = {]2, 3[}$$
	
	\section{Gráficos e Comportamento da Função}
	
	O traçado e o comportamento da curva $y = \log_a x$ dependem diretamente do valor da base $a$. Distinguem-se dois casos essenciais: a base maior do que um ($a > 1$) e a base entre zero e um ($0 < a < 1$).
	
	\subsection{Caso 1: Base Maior que a Unidade ($a > 1$)}
	
	Tomemos como modelo a função $f(x) = \log_2 x$. Atribuindo potências de $2$ para a variável independente $x$, calculamos os valores correspondentes de $y$:
	
	\begin{center}
		\begin{tabular}{|c|c|c|c|c|c|c|c|}
			\hline
			$x$ & $\frac{1}{4}$ & $\frac{1}{2}$ & $1$ & $2$ & $4$ & $8$ \\
			\hline
			$y = \log_2 x$ & $-2$ & $-1$ & $0$ & $1$ & $2$ & $3$ \\
			\hline
		\end{tabular}
	\end{center}
	
	Observando a tabela e a curva associada:
	\begin{itemize}[label=$\circ$]
		\item À medida que $x$ cresce, $y$ também cresce de modo contínuo. A função é \textbf{estritamente crescente}:
		$$x_1 < x_2 \iff \log_a x_1 < \log_a x_2 \quad (\text{para } a > 1)$$
		\item Para $x = 1$, temos $y = 0$. A curva intercepta o eixo horizontal no ponto $(1, 0)$.
		\item Para $x > 1$, os valores de $y$ são estritamente positivos.
		\item Para $0 < x < 1$, os valores de $y$ são negativos.
		\item À medida que $x$ se aproxima de zero pela direita ($x \to 0^+$), a curva aproxima-se indefinidamente do eixo vertical ($y \to -\infty$), fazendo do eixo das ordenadas uma \textbf{assíntota vertical}.
	\end{itemize}
	
	\begin{figure}[htbp]
		\centering
		\begin{tikzpicture}
			\begin{axis}[
				axis lines = center,
				xlabel = {$x$},
				ylabel = {$y$},
				xmin = -0.5, xmax = 8.5,
				ymin = -3.2, ymax = 3.5,
				domain = 0.12:8.2,
				samples = 150,
				grid = major,
				grid style = {dashed, gray!30},
				width = 9.5cm,
				height = 7.5cm,
				legend style = {at={(0.95,0.15)}, anchor=south east, font=\small}
				]
				\addplot[blue, thick] {ln(x)/ln(2)};
				\addlegendentry{$f(x) = \log_2 x$}
				
				% Ponto de interseção com o eixo x
				\addplot[mark=*, mark size=2pt, fill=black] coordinates {(1, 0)};
				\node[below right, font=\small] at (axis cs: 1, 0) {$(1, 0)$};
				
				% Ponto notável (2, 1)
				\addplot[mark=*, mark size=2pt, fill=black] coordinates {(2, 1)};
				\node[above left, font=\small] at (axis cs: 2, 1) {$(2, 1)$};
			\end{axis}
		\end{tikzpicture}
		\caption{Gráfico da função logarítmica crescente ($f(x) = \log_2 x$, com base $a = 2 > 1$).}
		\label{fig:log_crescente}
	\end{figure}
	
	\subsection{Caso 2: Base entre Zero e a Unidade ($0 < a < 1$)}
	
	Analisemos agora a função $g(x) = \log_{1/2} x$. Construindo a tabela para valores convenientes:
	
	\begin{center}
		\begin{tabular}{|c|c|c|c|c|c|c|c|}
			\hline
			$x$ & $\frac{1}{4}$ & $\frac{1}{2}$ & $1$ & $2$ & $4$ & $8$ \\
			\hline
			$y = \log_{1/2} x$ & $2$ & $1$ & $0$ & $-1$ & $-2$ & $-3$ \\
			\hline
		\end{tabular}
	\end{center}
	
	Neste caso:
	\begin{itemize}[label=$\circ$]
		\item À medida que $x$ aumenta, os valores de $y$ diminuem continuamente. A função é \textbf{estritamente decrescente}:
		$$x_1 < x_2 \iff \log_a x_1 > \log_a x_2 \quad (\text{para } 0 < a < 1)$$
		\item A curva mantém a interseção com o eixo das abscissas no ponto $(1, 0)$.
		\item Para $0 < x < 1$, os valores de $y$ são positivos, tendendo a $+\infty$ quando $x \to 0^+$.
		\item Para $x > 1$, os valores de $y$ tornam-se negativos.
	\end{itemize}
	
	\begin{figure}[htbp]
		\centering
		\begin{tikzpicture}
			\begin{axis}[
				axis lines = center,
				xlabel = {$x$},
				ylabel = {$y$},
				xmin = -0.5, xmax = 8.5,
				ymin = -3.5, ymax = 3.2,
				domain = 0.12:8.2,
				samples = 150,
				grid = major,
				grid style = {dashed, gray!30},
				width = 9.5cm,
				height = 7.5cm,
				legend style = {at={(0.95,0.85)}, anchor=north east, font=\small}
				]
				\addplot[red!80!black, thick] {ln(x)/ln(0.5)};
				\addlegendentry{$g(x) = \log_{1/2} x$}
				
				% Ponto de interseção com o eixo x
				\addplot[mark=*, mark size=2pt, fill=black] coordinates {(1, 0)};
				\node[above right, font=\small] at (axis cs: 1, 0) {$(1, 0)$};
				
				% Ponto notável (2, -1)
				\addplot[mark=*, mark size=2pt, fill=black] coordinates {(2, -1)};
				\node[below left, font=\small] at (axis cs: 2, -1) {$(2, -1)$};
			\end{axis}
		\end{tikzpicture}
		\caption{Gráfico da função logarítmica decrescente ($g(x) = \log_{1/2} x$, com base $a = \frac{1}{2}$).}
		\label{fig:log_decrescente}
	\end{figure}
	
	\section{A Relação de Inversão com a Função Exponencial}
	
	Uma propriedade geométrica e analítica de grande valor formativo é o vínculo de bijeção e inversão entre a função exponencial $e_a(x) = a^x$ e a função logarítmica $l_a(x) = \log_a x$.
	
	\begin{destaque}[Simetria em Relação à Bissetriz dos Quadrantes Ímpares]
		Como a função logarítmica é a inversa da exponencial correspondente:
		$$\log_a(a^x) = x \quad (\forall x \in \mathbb{R}) \quad \text{e} \quad a^{\log_a x} = x \quad (\forall x \in \mathbb{R}^*_+)$$
		o gráfico cartesiano de $y = \log_a x$ é a imagem refletida do gráfico de $y = a^x$ em relação à reta bissetriz dos quadrantes ímpares ($y = x$).
		
		A Figura \ref{fig:simetria_exp_log} apresenta essa simetria quando $a>0$ e a Figura  \ref{fig:simetria_exp_log_decrescente} quando $0<a<1$.
	\end{destaque}
		
	\begin{figure}[H]
		\centering
		\begin{tikzpicture}
			\begin{axis}[
				axis lines = center,
				xlabel = {$x$},
				ylabel = {$y$},
				xmin = -3.2, xmax = 5.5,
				ymin = -3.2, ymax = 5.5,
				grid = major,
				grid style = {dashed, gray!30},
				width = 10.5cm,
				height = 10.5cm,
				legend style = {
					at = {(0.5, -0.15)},
					anchor = north,
					legend columns = -1,
					font = \small
				},
				legend cell align = {left}
				]
				% Bissetriz dos quadrantes ímpares: y = x
				\addplot[gray!70, thick, dashdotted, domain=-3:5] {x};
				\addlegendentry{$y = x$}
				
				% Função Exponencial: f(x) = 2^x
				\addplot[blue, thick, domain=-3:2.4, samples=120] {2^x};
				\addlegendentry{$f(x) = 2^x$}
				
				% Função Logarítmica: g(x) = log_2(x)
				\addplot[red!80!black, thick, domain=0.12:5.2, samples=150] {ln(x)/ln(2)};
				\addlegendentry{$g(x) = \log_2 x$}
				
				% Segmentos conectando pares simétricos refletidos por y = x
				\draw[dashed, black!60] (axis cs: 0, 1) -- (axis cs: 1, 0);
				\draw[dashed, black!60] (axis cs: 1, 2) -- (axis cs: 2, 1);
				\draw[dashed, black!60] (axis cs: 2, 4) -- (axis cs: 4, 2);
				
				% Pontos notáveis na exponencial
				\addplot[only marks, mark=*, mark size=2pt, draw=blue, fill=blue] coordinates {(0, 1) (1, 2) (2, 4)};
				\node[left, font=\footnotesize, blue] at (axis cs: 0, 1.3) {$(0, 1)$};
				\node[left, font=\footnotesize, blue] at (axis cs: 1, 2) {$(1, 2)$};
				\node[above left, font=\footnotesize, blue] at (axis cs: 2, 4) {$(2, 4)$};
				
				% Pontos notáveis no logaritmo
				\addplot[only marks, mark=*, mark size=2pt, draw=red!80!black, fill=red!80!black] coordinates {(1, 0) (2, 1) (4, 2)};
				\node[below, font=\footnotesize, red!80!black] at (axis cs: 1.5, 0) {$(1, 0)$};
				\node[below right, font=\footnotesize, red!80!black] at (axis cs: 2, 1) {$(2, 1)$};
				\node[below right, font=\footnotesize, red!80!black] at (axis cs: 4, 2) {$(4, 2)$};
				
			\end{axis}
		\end{tikzpicture}
		
		\caption{Simetria entre as curvas de $f(x) = 2^x$ e $g(x) = \log_2 x$ em relação à bissetriz $y = x$.}
		\label{fig:simetria_exp_log}
	\end{figure}
	
	\begin{figure}[H]
		\centering
		\begin{tikzpicture}
			\begin{axis}[
				axis lines = center,
				xlabel = {$x$},
				ylabel = {$y$},
				xmin = -3.2, xmax = 5.5,
				ymin = -3.2, ymax = 5.5,
				grid = major,
				grid style = {dashed, gray!30},
				width = 10.5cm,
				height = 10.5cm,
				legend style = {
					at = {(0.5, -0.15)},
					anchor = north,
					legend columns = -1,
					font = \small
				},
				legend cell align = {left}
				]
				% Bissetriz dos quadrantes ímpares: y = x
				\addplot[gray!70, thick, dashdotted, domain=-3:5] {x};
				\addlegendentry{$y = x$}
				
				% Função Exponencial decrescente: f(x) = (1/2)^x
				\addplot[blue, thick, domain=-2.3:4.8, samples=120] {0.5^x};
				\addlegendentry{$f(x) = \left(\frac{1}{2}\right)^x$}
				
				% Função Logarítmica decrescente: g(x) = log_{1/2}(x)
				\addplot[red!80!black, thick, domain=0.12:4.8, samples=150] {ln(x)/ln(0.5)};
				\addlegendentry{$g(x) = \log_{1/2} x$}
				
				% Segmentos conectando pares ordenados simétricos (x, y) e (y, x)
				\draw[dashed, black!60] (axis cs: -1, 2) -- (axis cs: 2, -1);
				\draw[dashed, black!60] (axis cs: -2, 4) -- (axis cs: 4, -2);
				\draw[dashed, black!60] (axis cs: 0, 1) -- (axis cs: 1, 0);
				
				% Ponto comum de interseção na bissetriz (aproximadamente (0.641, 0.641))
				\addplot[mark=*, mark size=2.5pt, fill=black] coordinates {(0.641186, 0.641186)};
				
				% Pontos notáveis na exponencial
				\addplot[only marks, mark=*, mark size=2pt, draw=blue, fill=blue] coordinates {(-2, 4) (-1, 2) (0, 1) (1, 0.5)};
				\node[left, font=\footnotesize, blue] at (axis cs: -2, 4) {$(-2, 4)$};
				\node[left, font=\footnotesize, blue] at (axis cs: -1, 2) {$(-1, 2)$};
				\node[below left, font=\footnotesize, blue] at (axis cs: 0, 1) {$(0, 1)$};
				
				% Pontos notáveis no logaritmo
				\addplot[only marks, mark=*, mark size=2pt, draw=red!80!black, fill=red!80!black] coordinates {(4, -2) (2, -1) (1, 0) (0.5, 1)};
				\node[below, font=\footnotesize, red!80!black] at (axis cs: 4, -2) {$(4, -2)$};
				\node[below left, font=\footnotesize, red!80!black] at (axis cs: 2, -1) {$(2, -1)$};
				\node[below left, font=\footnotesize, red!80!black] at (axis cs: 1, 0) {$(1, 0)$};
				
			\end{axis}
		\end{tikzpicture}
		\caption{Simetria entre as curvas de $f(x) = \left(\frac{1}{2}\right)^x$ e $g(x) = \log_{1/2} x$ em relação à bissetriz $y = x$.}
		\label{fig:simetria_exp_log_decrescente}
	\end{figure}
	
	\section{Exercícios sobre a Função Logarítmica}
	
	\subsection*{Exercícios Resolvidos}
	
	\textbf{Exercício 1.} Determine o domínio da função real definida por:
	$$f(x) = \log_{(x - 1)} (x^2 - 4x - 5)$$
	
	\textit{Resolução:}  
	Para que a função esteja bem definida no conjunto dos números reais, é necessário satisfazer simultaneamente as três condições de existência:
	\begin{enumerate}[label=\textbf{\arabic*.}]
		\item \textbf{Logaritmando estritamente positivo:}
		$$x^2 - 4x - 5 > 0$$
		As raízes da equação correspondente $x^2 - 4x - 5 = 0$ são:
		$$x = \frac{4 \pm \sqrt{16 - 4(1)(-5)}}{2} = \frac{4 \pm 6}{2} \implies x_1 = -1 \quad \text{e} \quad x_2 = 5$$
		Como a parábola tem concavidade voltada para cima, a expressão é positiva para:
		$$x < -1 \quad \text{ou} \quad x > 5 \quad (S_1 = {]-\infty, -1[} \cup {]5, +\infty[})$$
		
		\item \textbf{Base estritamente positiva:}
		$$x - 1 > 0 \implies x > 1 \quad (S_2 = {]1, +\infty[})$$
		
		\item \textbf{Base diferente de $1$:}
		$$x - 1 \neq 1 \implies x \neq 2 \quad (S_3 = \mathbb{R} - \{2\})$$
	\end{enumerate}
	
	Determinando a interseção $S = S_1 \cap S_2 \cap S_3$:
	$$({]-\infty, -1[} \cup {]5, +\infty[}) \;\cap\; {]1, +\infty[} \;\cap\; (\mathbb{R} - \{2\}) = {]5, +\infty[}$$
	Portanto:
	$$\text{Dom}(f) = \{x \in \mathbb{R} \mid x > 5\} = {]5, +\infty[}$$
	
	\vspace{0.3cm}
	
	\textbf{Exercício 2.}  
	Considere a função real $g(x) = \log_2 (x - 3) + 1$.
	\begin{enumerate}[label=\textbf{\alph*)}]
		\item Determine o domínio, a imagem e a equação da assíntota vertical do gráfico de $g$.
		\item Calcule as coordenadas do ponto de interseção da curva com o eixo das abscissas (raiz de $g$).
	\end{enumerate}
	
	\textit{Resolução:}
	\begin{enumerate}[label=\textbf{\alph*)}]
		\item A condição de existência exige:
		$$x - 3 > 0 \implies x > 3$$
		Logo, $\text{Dom}(g) = {]3, +\infty[}$. O conjunto imagem de toda função logarítmica com translações verticais é $\text{Im}(g) = \mathbb{R}$.  
		A assíntota vertical da função elementar $y = \log_2 x$ (a reta $x = 0$) é transladada três unidades para a direita, resultando na reta de equação:
		$$x = 3$$
		
		\item A interseção com o eixo horizontal ocorre quando $g(x) = 0$:
		$$\log_2 (x - 3) + 1 = 0 \implies \log_2 (x - 3) = -1$$
		Pela definição fundamental de logaritmo:
		$$x - 3 = 2^{-1} \implies x - 3 = \frac{1}{2} \implies x = 3 + 0{,}5 = 3{,}5$$
		Assim, o gráfico intercepta o eixo $x$ no par ordenado $(3{,}5; 0)$.
	\end{enumerate}
	
	\vspace{0.3cm}
	
	\textbf{Exercício 3.}  
	Em fósseis e tecidos vegetais antigos, a idade $t$ (em anos) da amostra pode ser modelada a partir da fração restante $P$ de Carbono-14 em relação à sua quantidade na matéria viva, por meio da função logarítmica:
	$$t(P) = -8\,267 \cdot \ln P$$
	em que $P \in {]0, 1]}$ representa a razão $\frac{C(t)}{C_0}$ e $\ln$ denota o logaritmo natural.
	\begin{enumerate}[label=\textbf{\alph*)}]
		\item Justifique por que, embora a constante multiplicativa seja negativa ($-8\,267$), a idade $t(P)$ obtida é sempre não negativa no domínio considerado.
		\item Uma peça arqueológica de madeira encontrada em um sítio funerário revelou conter apenas $25\%$ ($P = 0{,}25$) da sua concentração original de Carbono-14. Determine a idade aproximada do fóssil (considere $\ln 2 \approx 0{,}693$).
	\end{enumerate}
	
	\textit{Resolução:}
	\begin{enumerate}[label=\textbf{\alph*)}]
		\item Para todo $P$ pertencente ao intervalo ${]0, 1]}$, o logaritmo natural é menor ou igual a zero ($\ln P \leqslant 0$), uma vez que a base $e \approx 2{,}718 > 1$ e $\ln 1 = 0$. Como o produto de dois valores não positivos é positivo, o fator $-8\,267$ assegura que $t(P) \geqslant 0$, em plena coerência com a contagem cronológica do tempo.
		
		\item Para $P = 0{,}25 = \frac{1}{4} = 2^{-2}$, reescrevemos o logaritmo:
		$$\ln(0{,}25) = \ln(2^{-2}) = -2 \cdot \ln 2$$
		Substituindo na expressão do tempo:
		$$t = -8\,267 \cdot (-2 \cdot \ln 2) = 16\,534 \cdot \ln 2$$
		Utilizando a aproximação $\ln 2 \approx 0{,}693$:
		$$t \approx 16\,534 \cdot 0{,}693 \approx 11\,458\text{ anos}$$
		Portanto, a madeira fóssil tem aproximadamente $11\,458$ anos de idade.
	\end{enumerate}
	
	\subsection*{Exercícios Propostos}
	
	\begin{enumerate}
		\item Determine o domínio das seguintes funções reais:
		\begin{multicols}{2}
			\begin{enumerate}[label=\textbf{\alph*)}]
				\item $f(x) = \log_3 (7 - 2x)$
				\item $g(x) = \log_5 (x^2 - 6x + 8)$
				\item $h(x) = \log_{(x - 3)} 10$
				\item $p(x) = \log_{(2x - 4)} (25 - x^2)$
			\end{enumerate}
		\end{multicols}
		
		\item Seja $f(x) = \log_a x$ uma função logarítmica cujo gráfico contém o ponto $A(8, -3)$.
		\begin{enumerate}[label=\textbf{\alph*)}]
			\item Determine a base $a$ da função.
			\item Calcule o valor de $f(32)$ e de $f\left(\frac{1}{16}\right)$.
			\item A função $f$ é crescente ou decrescente? Justifique a resposta com base no valor de $a$.
		\end{enumerate}
		
		\item Para cada uma das funções a seguir, determine a raiz, o domínio e a equação da assíntota vertical do respectivo gráfico:
		\begin{enumerate}[label=\textbf{\alph*)}]
			\item $f(x) = \log_3 (x + 2)$
			\item $g(x) = 2 - \log_2 x$
			\item $h(x) = \log_{1/3} (2x - 5) + 2$
		\end{enumerate}
		
		\item A altitude $h$ (em quilômetros acima do nível do mar) pode ser estimada a partir da pressão atmosférica local $P$ (em milímetros de mercúrio, $\text{mmHg}$) através do modelo logarítmico:
		$$h(P) = -5{,}6 \cdot \ln\left(\frac{P}{760}\right)$$
		em que $760\text{ mmHg}$ corresponde à pressão média padrão ao nível do mar ($h = 0$).
		\begin{enumerate}[label=\textbf{\alph*)}]
			\item Mostre que, para qualquer pressão inferior à do nível do mar ($0 < P < 760$), a altitude estimada $h(P)$ é estritamente positiva.
			\item Um grupo de alpinistas mediu uma pressão de $380\text{ mmHg}$ no acampamento alto de uma montanha. Determine a altitude em que se encontram (considere $\ln 2 \approx 0{,}693$).
		\end{enumerate}
	\end{enumerate}
	
	\chapter{Equações Exponenciais e Logarítmicas}
	\label{chap:equacoes_inequacoes}
	
	Neste capítulo, estudamos os métodos de resolução para equações exponenciais e logarítmicas. Investigamos desde relações exponenciais que não admitem redução direta a bases comuns até equações em que a incógnita figura no logaritmando ou na base do logaritmo, momento em que as condições de existência orientam a validação das soluções encontradas.
	
	\section{Equações Exponenciais Não Redutíveis a Bases Comuns}
	
	No estudo inicial das equações exponenciais, os métodos baseiam-se quase exclusivamente em manipulações algébricas elementares para igualar as bases de ambos os membros da igualdade:
	$$a^{f(x)} = a^{g(x)} \implies f(x) = g(x)$$
	
	Entretanto, em inúmeras situações matemáticas e em modelagens de fenômenos contínuos, deparamo-nos com equações em que os dois membros envolvem fatores primos distintos que não admitem redução direta a uma base comum, tais como:
	$$a^x = b \quad \text{ou} \quad a^{f(x)} = b^{g(x)} \quad (a \neq b)$$
	
	Nesses casos, o método algébrico apropriado consiste na aplicação de logaritmos em ambos os membros da igualdade, permitindo transferir os expoentes para a linha multiplicativa por meio da propriedade da potência.
	
	\subsection{Procedimento Geral de Resolução}
	
	Dada a equação geral:
	$$a^{f(x)} = b^{g(x)}$$
	com $a, b \in \mathbb{R}^*_+$, $a \neq 1$ e $b \neq 1$, o procedimento analítico organiza-se nas seguintes etapas:
	
	\begin{destaque}[Roteiro de Resolução via Logaritmos]
		\begin{enumerate}[label=\textbf{\arabic*.}]
			\item \textbf{Aplicação do logaritmo:}  
			Aplica-se o logaritmo em uma base conveniente (geralmente decimal ou natural, ou na própria base de uma das potências) em ambos os lados da igualdade:
			$$\log(a^{f(x)}) = \log(b^{g(x)})$$
			
			\item \textbf{Descida dos expoentes:}  
			Pela propriedade do logaritmo de uma potência:
			$$f(x) \cdot \log a = g(x) \cdot \log b$$
			
			\item \textbf{Isolamento da incógnita:}  
			A equação transforma-se em uma equação algébrica linear ou polinomial na variável $x$, permitindo isolá-la diretamente em termos dos logaritmos das constantes envolvidas.
		\end{enumerate}
	\end{destaque}
	
	\subsection*{Exemplos}
	
	\begin{enumerate}[label=\textbf{\alph*)}]
		\item \textbf{Equação com potências simples:}  
		Determine a solução exata e a aproximação decimal de $2^x = 5$, considerando $\log 2 \approx 0{,}301$ e $\log 5 \approx 0{,}699$:
		$$\log(2^x) = \log 5 \implies x \cdot \log 2 = \log 5 \implies x = \frac{\log 5}{\log 2} = \log_2 5$$
		Substituindo as aproximações:
		$$x \approx \frac{0{,}699}{0{,}301} \approx 2{,}322$$
		
		\item \textbf{Equação com expoentes lineares em ambas as bases:}  
		Resolva $3^{2x - 1} = 5^{x + 2}$:
		\begin{align*}
			\log(3^{2x - 1}) &= \log(5^{x + 2}) \\
			(2x - 1) \cdot \log 3 &= (x + 2) \cdot \log 5 \\
			2x \log 3 - \log 3 &= x \log 5 + 2 \log 5 \\
			2x \log 3 - x \log 5 &= 2 \log 5 + \log 3 \\
			x(2\log 3 - \log 5) &= 2 \log 5 + \log 3 \\
			x &= \frac{2\log 5 + \log 3}{2\log 3 - \log 5} = \frac{\log(25 \cdot 3)}{\log(9 / 5)} = \frac{\log 75}{\log 1{,}8}
		\end{align*}
	\end{enumerate}
	
	\subsection*{Exercícios Resolvidos}
	
	\textbf{Exercício 1.} Resolva a equação $7^x = 12$, expressando a solução em forma exata e em valor aproximado (adote $\log 2 \approx 0{,}30$, $\log 3 \approx 0{,}48$ e $\log 7 \approx 0{,}84$).
	
	\textit{Resolução:}  
	Aplicamos o logaritmo decimal em ambos os membros:
	$$\log(7^x) = \log 12 \implies x \cdot \log 7 = \log 12 \implies x = \frac{\log 12}{\log 7}$$
	Decompomos o numerador em fatores primos ($12 = 2^2 \cdot 3$):
	$$\log 12 = \log(2^2 \cdot 3) = 2\log 2 + \log 3 \approx 2 \cdot (0{,}30) + 0{,}48 = 1{,}08$$
	Substituindo os valores:
	$$x \approx \frac{1{,}08}{0{,}84} \approx 1{,}286$$
	Portanto, $S = \left\{\frac{\log 12}{\log 7}\right\} \approx \{1{,}286\}$.
	
	\vspace{0.3cm}
	
	\textbf{Exercício 2.} Resolva a equação $2^{x+1} \cdot 5^{2x} = 30$, admitindo $\log 2 \approx 0{,}30$ e $\log 3 \approx 0{,}48$.
	
	\textit{Resolução:}  
	Aplicamos o logaritmo decimal à igualdade:
	$$\log(2^{x+1} \cdot 5^{2x}) = \log 30$$
	Aplicando as propriedades do produto e da potência:
	$$\log(2^{x+1}) + \log(5^{2x}) = \log(3 \cdot 10)$$
	$$(x + 1) \cdot \log 2 + 2x \cdot \log 5 = \log 3 + \log 10$$
	Lembrando que $\log 5 = \log\left(\frac{10}{2}\right) = \log 10 - \log 2 = 1 - \log 2$:
	$$x \log 2 + \log 2 + 2x(1 - \log 2) = \log 3 + 1$$
	$$x \log 2 + \log 2 + 2x - 2x \log 2 = \log 3 + 1$$
	$$x(2 - \log 2) = \log 3 + 1 - \log 2$$
	$$x=\frac{\log 3 + 1 - \log 2}{2 - \log 2}$$
	
	Substituindo as aproximações $\log 2 \approx 0{,}30$ e $\log 3 \approx 0{,}48$:
	$$x = \frac{0{,}48 + 1 - 0{,}30}{2 - 0{,}30}$$
	$$x = \frac{1{,}18}{1{,}70} \approx 0{,}694$$
	Portanto, $S \approx \{0{,}694\}$.
	
	\vspace{0.3cm}
	
	\textbf{Exercício 3.} Em um experimento biológico, uma população bacteriana $P_1(t) = 800 \cdot (1{,}4)^t$ compete com outra cultura cujo crescimento segue $P_2(t) = 200 \cdot (2{,}1)^t$, onde $t$ é medido em dias. Determine após quantos dias as duas populações atingirão exatamente o mesmo número de indivíduos (adote $\log 2 \approx 0{,}30$, $\log 3 \approx 0{,}48$ e $\log 7 \approx 0{,}84$).
	
	\textit{Resolução:}  
	Igualando as duas populações:
	$$800 \cdot (1{,}4)^t = 200 \cdot (2{,}1)^t$$
	Dividindo ambos os membros por $200$:
	$$4 \cdot (1{,}4)^t = (2{,}1)^t \implies 4 = \frac{(2{,}1)^t}{(1{,}4)^t} \implies 4 = \left(\frac{2{,}1}{1{,}4}\right)^t$$
	Simplificando a fração decimal:
	$$\frac{2{,}1}{1{,}4} = \frac{21}{14} = \frac{3}{2} = 1{,}5$$
	A equação assume a forma:
	$$\left(\frac{3}{2}\right)^t = 4$$
	Aplicando o logaritmo decimal em ambos os lados:
	$$\log\left(\left(\frac{3}{2}\right)^t\right) = \log 4 \implies t \cdot \log\left(\frac{3}{2}\right) = \log(2^2)$$
	$$t \cdot (\log 3 - \log 2) = 2 \cdot \log 2$$
	Isolando $t$:
	$$t = \frac{2 \cdot \log 2}{\log 3 - \log 2}$$
	Substituindo os dados numéricos:
	$$t \approx \frac{2 \cdot 0{,}30}{0{,}48 - 0{,}30} = \frac{0{,}60}{0{,}18} = \frac{60}{18} = \frac{10}{3} \approx 3{,}33\text{ dias}$$
	Portanto, as culturas atingem o mesmo número de bactérias após $3\text{ dias e } 8\text{ horas}$.
	
	\subsection*{Exercícios Propostos}
	
	\begin{enumerate}
		\item Resolva as seguintes equações exponenciais em $\mathbb{R}$, determinando a solução exata em termos de logaritmos decimais:
		\begin{enumerate}[label=\textbf{\alph*)}]
			\item $3^x = 20$
			\item $5^{2x - 1} = 7$
			\item $2^{x + 3} = 3^{x - 1}$
			\item $4^x \cdot 3^{x + 1} = 5^x$
		\end{enumerate}
		
		\item Considerando os valores aproximados $\log 2 \approx 0{,}301$ e $\log 3 \approx 0{,}477$, determine a solução numérica aproximada de:
		\begin{enumerate}[label=\textbf{\alph*)}]
			\item $6^x = 10$
			\item $3^{x + 2} = 2^{2x + 1}$
		\end{enumerate}
		
		\item Um equipamento eletrônico sofre desvalorização exponencial com o passar do tempo, de modo que o seu valor de mercado $V(t)$ (em reais) após $t$ anos de uso é modelado por $V(t) = 6\,000 \cdot (0{,}82)^t$. Determine em quantos anos o valor do equipamento se reduzirá à metade do valor original de aquisição (adote $\log 2 \approx 0{,}30$ e $\log 8{,}2 \approx 0{,}91$).
		
		\item Resolva a equação $e^{2x} - 4e^x + 3 = 0$, determinando suas raízes em termos de logaritmos naturais ($\ln$).
	\end{enumerate}
	
	\section{Equações Logarítmicas}
	
	Uma equação é dita logarítmica quando a incógnita aparece vinculada a pelo menos um operador logarítmico. Classificamos os procedimentos de resolução em três tipos fundamentais.
	
	\subsection{1º Tipo: Redução a uma Igualdade de Mesma Base}
	
	São as equações que, após simplificações operatórias, assumem a forma:
	$$\log_a f(x) = \log_a g(x)$$
	com $a > 0$ e $a \neq 1$.
	
	Pela propriedade da unicidade dos logaritmos, essa igualdade equivale à igualdade dos logaritmandos, desde que respeitadas as condições de existência:
	$$\begin{cases}
		f(x) = g(x) \\
		f(x) > 0 \\
		g(x) > 0
	\end{cases}$$
	
	\newpage
	
	\begin{destaque}[Exemplo Resolvido]
		Resolva a equação:
		$$\log_2 (2x - 3) = \log_2 (x + 1)$$
		
		\textit{Resolução:}
		\begin{enumerate}[label=\textbf{\arabic*.}]
			\item Condições de existência:
			$$\begin{cases} 
				2x - 3 > 0 \implies x > \frac{3}{2} \\
				x + 1 > 0 \implies x > -1
			\end{cases} \implies x \in {]\frac{3}{2}, +\infty[}$$
			
			\item Igualando os logaritmandos:
			$$2x - 3 = x + 1 \implies x = 4$$
			
			\item Verificação:
			Como $4 \in {]\frac{3}{2}, +\infty[}$, o valor é admitido.
		\end{enumerate}
		Portanto, o conjunto solução é $S = \{4\}$.
	\end{destaque}
	
	\subsection{2º Tipo: Aplicação Direta da Definição}
	
	São as equações que se reduzem à forma elementar:
	$$\log_a f(x) = k \quad (k \in \mathbb{R})$$
	cuja resolução decorre imediatamente da definição fundamental:
	$$f(x) = a^k \quad \text{com } f(x) > 0$$
	
	\begin{destaque}[Exemplo Resolvido]
		Resolva a equação:
		$$\log_3 (x^2 - 7) = 2$$
		
		\textit{Resolução:}
		\begin{enumerate}[label=\textbf{\arabic*.}]
			\item Condição de existência:
			$$x^2 - 7 > 0 \implies x < -\sqrt{7} \quad \text{ou} \quad x > \sqrt{7}$$
			
			\item Pela definição de logaritmo:
			$$x^2 - 7 = 3^2 \implies x^2 - 7 = 9 \implies x^2 = 16 \implies x = \pm 4$$
			
			\item Verificação:
			Como $(\pm 4)^2 = 16 > 7$, ambos os valores satisfazem a condição de existência.
		\end{enumerate}
		Portanto, o conjunto solução é $S = \{-4, 4\}$.
	\end{destaque}
	
	\subsection{3º Tipo: Resolução por Mudança de Variável}
	
	Ocorrem quando a expressão envolve potências do próprio logaritmo. Nesses casos, introduzimos uma incógnita auxiliar $y = \log_a x$, transformando a expressão em uma equação polinomial.
	
	\begin{destaque}[Exemplo Resolvido]
		Resolva a equação:
		$$(\log_2 x)^2 - 3\log_2 x + 2 = 0$$
		
		\textit{Resolução:}
		\begin{enumerate}[label=\textbf{\arabic*.}]
			\item Condição de existência:
			$$x > 0$$
			
			\item Mudança de variável:
			Fazendo $y = \log_2 x$, obtemos:
			$$y^2 - 3y + 2 = 0 \implies (y - 1)(y - 2) = 0 \implies y = 1 \quad \text{ou} \quad y = 2$$
			
			\item Retorno à variável original:
			\begin{itemize}[label=$\circ$]
				\item Para $y = 1$: $\log_2 x = 1 \implies x = 2^1 = 2$.
				\item Para $y = 2$: $\log_2 x = 2 \implies x = 2^2 = 4$.
			\end{itemize}
			Ambos atendem à restrição $x > 0$.
		\end{enumerate}
		Portanto, o conjunto solução é $S = \{2, 4\}$.
	\end{destaque}
	
	\section{Exercícios Resolvidos e Propostos}
	
	\subsection*{Exercícios Resolvidos}
	
	\textbf{Exercício 1.} Resolva a equação:
	$$\log (x + 2) + \log (x - 1) = 1$$
	
	\textit{Resolução:}
	\begin{enumerate}[label=\textbf{\arabic*.}]
		\item Condições de existência:
		$$\begin{cases}
			x + 2 > 0 \implies x > -2 \\
			x - 1 > 0 \implies x > 1
		\end{cases} \implies x \in {]1, +\infty[}$$
		
		\item Aplicando a propriedade do logaritmo do produto:
		$$\log[(x + 2)(x - 1)] = 1$$
		Como a base implícita é $10$, pela definição fundamental:
		$$(x + 2)(x - 1) = 10^1 \implies x^2 + x - 2 = 10 \implies x^2 + x - 12 = 0$$
		
		\item Determinando as raízes:
		$$(x + 4)(x - 3) = 0 \implies x = -4 \quad \text{ou} \quad x = 3$$
		
		\item Análise de pertinência ao domínio:
		\begin{itemize}[label=$\circ$]
			\item $x = -4 \notin {]1, +\infty[}$ (solução estranha decorrente do produto de dois fatores negativos).
			\item $x = 3 \in {]1, +\infty[}$.
		\end{itemize}
	\end{enumerate}
	Portanto, o conjunto solução é $S = \{3\}$.
	
	\vspace{0.3cm}
	
	\textbf{Exercício 2.} Um investimento financeiro de capital inicial $C_0$ rende juros compostos contínuos segundo o modelo de montante $M(t) = C_0 \cdot (1{,}05)^t$, onde $t$ representa o tempo em anos. Deseja-se determinar em quanto tempo esse investimento pelo menos triplica o capital aplicado.
	\begin{enumerate}[label=\textbf{\alph*)}]
		\item Modele o problema por meio de uma inequação logarítmica.
		\item Calcule o tempo mínimo inteiro de anos necessário, considerando $\log 3 \approx 0{,}477$ e $\log 1{,}05 \approx 0{,}021$.
	\end{enumerate}
	
	\textit{Resolução:}
	\begin{enumerate}[label=\textbf{\alph*)}]
		\item Para que o montante triplique:
		$$C_0 \cdot (1{,}05)^t \geqslant 3 C_0 \implies (1{,}05)^t \geqslant 3$$
		Como a base decimal é maior que 1, aplicamos o logaritmo mantendo o sentido:
		$$\log(1{,}05^t) \geqslant \log 3 \implies t \cdot \log 1{,}05 \geqslant \log 3 \implies t \geqslant \frac{\log 3}{\log 1{,}05}$$
		
		\item Substituindo as aproximações:
		$$t \geqslant \frac{0{,}477}{0{,}021} \approx 22{,}71\text{ anos}$$
		Como se busca o menor número inteiro de anos, temos $t = 23\text{ anos}$.
	\end{enumerate}
	
	\newpage
	
	\subsection*{Exercícios Propostos}
	
	\begin{enumerate}
		\item Resolva as seguintes equações logarítmicas em $\mathbb{R}$:
		\begin{enumerate}[label=\textbf{\alph*)}]
			\item $\log_3 (4x - 5) = 2$
			\item $\log_2 (x^2 - 5x + 8) = 1$
			\item $\log_4 x + \log_4 (x - 6) = 2$
			\item $(\log_3 x)^2 - 4\log_3 x + 3 = 0$
		\end{enumerate}
		
		\item Um reservatório de resíduos industriais sofre vazamento acidental. A concentração de poluentes na água $C(t)$, medida em partes por milhão ($\text{ppm}$), diminui conforme o tempo $t$ (em dias) segundo $C(t) = 400 \cdot (0{,}5)^{t/4}$. A legislação ambiental permite o uso da água para abastecimento apenas quando a concentração for menor que $25\text{ ppm}$.
		\begin{enumerate}[label=\textbf{\alph*)}]
			\item Escreva a inequação correspondente a essa exigência ambiental.
			\item Determine a partir de quantos dias a água do reservatório estará liberada para captação.
		\end{enumerate}
	\end{enumerate}
\end{document}